% Quadratische Gleichungen – bank/quadratische-gleichungen/ (zusammenbau v0.4, Heft)
\einheitenkopf[t3]{Quadratische Gleichungen}
\setcounter{aufgabe}{18}

% Anlauf (ohne Prüfkennung)
\begin{aufgabe}{Wurzelziehen und Lösbarkeit – Anlauf}
\begin{teile}
\teil $(x - 3)^2 = 25$ – wie viele Lösungen?\\ \kreuz{zwei} \\ \kreuz{eine} \\ \kreuz{keine}
\end{teile}
\begin{gleichungsraster}[2]
\gl{x^2 = 81} &
\end{gleichungsraster}
\begin{teile}
\teil $x^2 = -100$ – Lösungen? Begründe.
\end{teile}
\end{aufgabe}

\begin{aufgabe}{Wurzelziehen und Lösbarkeit – Prüfungsaufgaben}
\begin{teile}
\steil Wahr oder falsch? Gib außerdem eine quadratische Gleichung ohne Lösung an. \hfill (P10 2021 OS)\\ A: $(x - 6)^2 = 0$ hat genau eine Lösung. \kreuz{wahr} \kreuz{falsch}\\ B: $7$ und $11$ sind Lösungen von $(x + 9)^2 = 4$. \kreuz{wahr} \kreuz{falsch} Gleichung ohne Lösung: \leerfeld
\steil Wahr oder falsch? Gib außerdem eine quadratische Gleichung ohne Lösung an. \hfill (P10 2021 OS)\\ A: $(x + 10)^2 = 0$ hat keine Lösung. \kreuz{wahr} \kreuz{falsch}\\ B: $-1$ und $-11$ sind Lösungen von $(x + 6)^2 = 25$. \kreuz{wahr} \kreuz{falsch} Gleichung ohne Lösung: \leerfeld
\end{teile}
\end{aufgabe}

% Anlauf (ohne Prüfkennung)
\begin{aufgabe}{Nullprodukt – Anlauf}
\begin{teile}
\teil $(x - 2) \cdot (x + 7) = 0$ – gilt der Satz vom Nullprodukt?\\ \kreuz{gilt} \\ \kreuz{gilt nicht}
\end{teile}
\begin{gleichungsraster}[2]
\gl{(x - 2) \cdot (x - 5) = 0} &
\gl{(x + 4) \cdot (x - 1) = 0} \\
\end{gleichungsraster}
\end{aufgabe}

\begin{aufgabe}{Nullprodukt – Prüfungsaufgaben}
\begin{teile}
\teil Welcher Wert erfüllt $x \cdot (x + 9) = -14$? \hfill (P10 2025 OS)\\ \kreuz{$x = 2$} \\ \kreuz{$x = 7$} \\ \kreuz{$x = -7$} \\ \kreuz{$x = -14$}
\teil Welcher Wert erfüllt $x \cdot (x + 7) = -10$? \hfill (P10 2025 OS)\\ \kreuz{$x = 2$} \\ \kreuz{$x = 5$} \\ \kreuz{$x = -2$} \\ \kreuz{$x = -10$}
\steil Nullstellen von $y = 3x^2 + 15x + 12$ durch Faktorisieren? Teile erst durch $3$ und suche zwei Zahlen mit dem Produkt $4$ und der Summe $5$. \hfill (P10 2020 OS) \\ x1 = \leerfeld , x2 = \leerfeld
\steil Nullstellen von $y = 2x^2 - 2x - 12$ durch Faktorisieren? Teile erst durch $2$ und suche zwei Zahlen mit dem Produkt $-6$ und der Summe $-1$. \hfill (P10 2020 OS) \\ x1 = \leerfeld , x2 = \leerfeld
\end{teile}
\end{aufgabe}

% Anlauf (ohne Prüfkennung)
\begin{aufgabe}{p-q-Formel – Anlauf}
\begin{teile}
\teil $x^2 - x - 12 = 0$ – p und q? p = \leerfeld , q = \leerfeld
\end{teile}
\begin{gleichungsraster}[2]
\gl{x^2 + 6x + 8 = 0} &
\gl{x^2 - 4x - 12 = 0} \\
\end{gleichungsraster}
\end{aufgabe}

\begin{aufgabe}{p-q-Formel – Prüfungsaufgaben}
\begin{teile}
\steil Die Funktion $f$ mit $f(x) = 2x^2 - 2x - 40$ – Nullstellen? \hfill (P10 2020 OS) \\ x1 = \leerfeld , x2 = \leerfeld
\steil Die Funktion $f$ mit $f(x) = 2x^2 + 16x + 30$ – Nullstellen? \hfill (P10 2020 OS) \\ x1 = \leerfeld , x2 = \leerfeld
\steil Die Funktion $g$ mit $g(x) = x^2 - 10x + 22$ – Schnittpunkte mit der x-Achse (zwei Stellen)? \hfill (P10 2025 OS) \\ ( \leerfeld | 0), ( \leerfeld | 0)
\steil Die Funktion $g$ mit $g(x) = x^2 + 4x - 1$ – Schnittpunkte mit der x-Achse (zwei Stellen)? \hfill (P10 2025 OS) \\ ( \leerfeld | 0), ( \leerfeld | 0)
\steil Parabel $y = -x^2 + 2x + 6$ – an welchen Stellen $x$ ist $y = -2$? \hfill (P10 2023 OS) \\ x1 = \leerfeld , x2 = \leerfeld
\steil Parabel $y = -x^2 - 6x + 4$ – an welchen Stellen $x$ ist $y = -12$? \hfill (P10 2023 OS) \\ x1 = \leerfeld , x2 = \leerfeld
\steil Parabel $y = x^2 + 3x - 2$ und Gerade $y = 2x + 4$ – Schnittpunkte? \hfill (P10 2021 OS) \\ ( \leerfeld | \leerfeld ), ( \leerfeld | \leerfeld )
\steil Parabel $y = x^2 + 4x - 3$ und Gerade $y = x + 1$ – Schnittpunkte? \hfill (P10 2021 OS) \\ ( \leerfeld | \leerfeld ), ( \leerfeld | \leerfeld )
\steil Parabel $y = (x - 3)^2 - 2$ und Gerade $y = -2x + 19$ – Schnittpunkte? \hfill (P10 2022 OS) \\ ( \leerfeld | \leerfeld ), ( \leerfeld | \leerfeld )
\steil Parabel $y = (x + 1)^2 - 3$ und Gerade $y = x + 10$ – Schnittpunkte? \hfill (P10 2022 OS) \\ ( \leerfeld | \leerfeld ), ( \leerfeld | \leerfeld )
\steil Parabel $y = x^2 - 2$ und Gerade $y = 3x + 2$ – Schnittpunkte? \hfill (P10 2024 OS) \\ ( \leerfeld | \leerfeld ), ( \leerfeld | \leerfeld )
\steil Parabel $y = x^2 - 8$ und Gerade $y = -3x + 2$ – Schnittpunkte? \hfill (P10 2024 OS) \\ ( \leerfeld | \leerfeld ), ( \leerfeld | \leerfeld )
\steil Zeige, dass $y = (x - 2)^2 - 3$ die Normalform $y = x^2 - 4x + 1$ hat, und berechne die Nullstellen (zwei Stellen). \hfill (P10 2017 OS) \\ x1 ≈ \leerfeld , x2 ≈ \leerfeld
\steil Zeige, dass $y = (x + 1)^2 - 5$ die Normalform $y = x^2 + 2x - 4$ hat, und berechne die Nullstellen (zwei Stellen). \hfill (P10 2017 OS) \\ x1 ≈ \leerfeld , x2 ≈ \leerfeld
\end{teile}
\end{aufgabe}
